\subsection{A universal obstruction for the two-circle chart family}

\begin{lemma}[A right-zero action has a natural graph model]
\label{C:right-zero-action-graph}
Let $R=\{1,p,q\}$, with $ab=b$ for $a,b\in\{p,q\}$,
act on the left on a finite set $T$.  Its two nonidentity
transformations $P,Q$ are retractions onto the same subset
$U\subseteq T$.  The action category has a natural classifying
space equivalence to the graph with vertices $U$ and, for each
$t\in T\setminus U$, one oriented edge from $P(t)$ to $Q(t)$.
An equivariant set map sends this edge to the edge indexed by
the image of $t$, or collapses it to a vertex if that image
belongs to the new common retraction image.
\end{lemma}
\begin{proof}
The equations $PQ=Q$ and $QP=P$ show that
$\operatorname{im}P=\operatorname{im}Q=U$, and both maps
are the identity on $U$.  In the action category every vertex
of $U$ has endomorphism monoid $R$.  A vertex outside $U$ has
only its identity and two outgoing arrows, labeled $p$ and $q$,
to $P(t)$ and $Q(t)$.  There are no other arrows between
different objects.  Postcomposition fixes each of these two
arrow labels, since $ap=p$ and $aq=q$ for $a\in R$.

Consequently its classifying space is the homotopy colimit of
a graph of spaces: the vertices in $U$ carry $BR$, outside
vertices carry points, and each of the two incident edges at
an outside vertex carries $BR$, mapping identically to its
$U$-endpoint and constantly to the outside point.  The space
$BR$ is contractible, since $R$ has a right zero.  Collapsing
each copy of $BR$ gives an equivalence to the indicated
bipartite graph.  Removing its bivalent outside vertices gives
the graph in the statement.

This equivalence can be chosen naturally, rather than merely
separately for each action.  A nerve simplex either stays at a
vertex of $U$, or has one initial arrow $p$ or $q$ from an
outside vertex, followed by endomorphisms of its endpoint.
Collapse the first kind to that vertex; send the second along
the corresponding half-edge, with position equal to the sum
of its endpoint barycentric coordinates.  The right-zero
identities make this compatible with every face and degeneracy.
An equivariant map preserves that half-edge label, unless it
sends its outside vertex into $U$, in which case both sides
collapse it.  This proves the asserted naturality.
\end{proof}

\begin{lemma}[Two commuting right-zero actions]
\label{C:two-retraction-cubical-model}
Suppose that $P_0,P_1$ and $Q_0,Q_1$ are two commuting pairs
of retractions on a finite set $T$, each pair having a common
image.  Put $U=\operatorname{im}P_0=\operatorname{im}P_1$
and $V=\operatorname{im}Q_0=\operatorname{im}Q_1$.
The double action correspondence has a natural finite CW model
$K(T)$ with cells
\[
\begin{array}{c|l}
0&U\cap V\\
1&U\setminus V\quad\text{and}\quad V\setminus U\\
2&T\setminus(U\cup V).
\end{array}
\]
Its two-cell indexed by $t$ has the four faces indexed by
$P_0t,P_1t,Q_0t,Q_1t$, with a face collapsed when its index
lies in $U\cap V$.  Equivariant maps preserve these cells
or collapse a coordinate.
\end{lemma}
\begin{proof}
First apply Lemma~\ref{C:right-zero-action-graph} to the $P$
action.  The $Q$ action commutes with it and therefore acts
on its graph by the natural cellular maps in that lemma.
Regard the graph as a one-dimensional simplicial set, including
its degenerate edges.  Its set of vertices is $U$, and its set
of edges together with degenerate edges is literally $T$.
Apply the same graph construction to the $Q$ action in each
simplicial degree.  Naturality makes this a simplicial graph.
Iterated realization and the two action bars commute, and
the levelwise graph equivalences give the asserted equivalence.
The degeneracies are subcomplex inclusions, so this use of
levelwise equivalences is valid.

In degree zero the graph has vertices $U\cap V$ and edges
$U\setminus V$.  In degree one, the new vertices are
$V\setminus U$ and the new edges are
$T\setminus(U\cup V)$.  Higher simplicial degrees are
degenerate copies of these.  Realization therefore gives
exactly the displayed finite two-dimensional cell structure.
The four boundary maps are the original four retractions;
their corner equations are precisely their commutation
relations.  This also proves the stated cellwise naturality.
\end{proof}

\begin{theorem}[The two-circle chart family has zero degree-two products]
\label{C:affine-chart-family-zero}
Let $C$ have objects $0,1,2$ and only forward cross arrows.
Let $\operatorname{End}(0)$ be the identity-adjoined two-element
left-zero band,
$\operatorname{End}(1)=1$, and let
$\operatorname{End}(2)$ be the identity-adjoined two-element
right-zero band.  Thus both these corner monoids have order three.
Suppose
\[
 C(0,1)=\{a_0,a_1\},\qquad C(1,2)=\{b_0,b_1\},
\]
with every nonidentity element of the first corner sending
both $a_i$ to $a_0$ by precomposition, and every nonidentity
element of the last corner sending both $b_i$ to $b_0$ by
postcomposition.  The set $C(0,2)$ may be any finite nonempty
biset, and composition may be any equivariant map.
Then every rational product of degree-two classes on $BC$ is
zero.  In particular this family never realizes
$\mathbb {CP}^2$, regardless of its total number of arrows.
\end{theorem}
\begin{proof}
Put $T=C(0,2)$ and $t=b_1a_1$.  Precomposition gives two
retractions $P_0,P_1$ with common image $U$, and postcomposition
gives two retractions $Q_0,Q_1$ with common image $V$.
They commute.  Equivariance says
\[
 P_0t=P_1t=b_1a_0,\qquad
 Q_0t=Q_1t=b_0a_1.
\]
These equations determine the entire composition map from $t$.

The bar decomposition by chains of distinct objects expresses
$BC$ as a triangle of spaces.  Its vertex spaces are the three
corner classifying spaces; all are contractible.  The two
adjacent edge correspondences have the graph models
$A=B=S^1$: each has one fixed state and one transient state
whose two retraction images agree.  The long edge has the
model $K=K(T)$ of Lemma~\ref{C:two-retraction-cubical-model}.
Because the middle corner is trivial, the face correspondence
is the product $P=A\times B=T^2$.  Its three face maps are
the two projections and the cellular map
\[
 h:A\times B\longrightarrow K
\]
induced by composition.  These comparisons are compatible
with the face maps by the naturality in the two lemmas.
Thus replacing the original bar diagram by this diagram is
an objectwise equivalence, and gives the same homotopy colimit.
Equivalently, the triangle of spaces is obtained from the
three edge mapping cylinders by attaching
$P\times\Delta^2$ along its three faces.

The product of the two one-edge circle models has exactly
one two-cell.  Its image under $h$ is the two-cell indexed
by $t$, if $t\notin U\cup V$, and is degenerate otherwise.
In the first case its opposite faces agree by the two
displayed equations.  Therefore this cell is a cellular
two-cycle.  It represents a nonzero primitive class of
$H_2(K;\mathbb Z)$: its coefficient on its own basis cell
is one, and $K$ has no three-cells.

For clarity, the homology spectral sequence of the triangle
filtration has, in total degree four, only
$H_2(P)$ in triangle degree two.  Its first differential is
the map to the sum of the second homology of the three edges,
namely
\[
 h_*:H_2(T^2;\mathbb Q)\longrightarrow H_2(K;\mathbb Q),
\]
since the other two edges are circles.  No higher differential
can hit a nonzero vertex group in this degree.  The preceding
cell calculation makes $h_*$ injective.  Thus
$H_4(BC;\mathbb Q)=0$, and all degree-two products vanish.

In the remaining case, suppose for example that $t\in U$.
Both $P_i$ fix $t$, so the composition map is independent of
the $a$-coordinate.  On the natural graph models,
$h$ factors through the projection to $B$:
$h=\alpha\operatorname{pr}_B$ for a cellular map
$\alpha:B\to K$.  There is then an actual map from the
triangle of spaces to the unreduced suspension $\Sigma K$.
On the long edge use the identity suspension, on the $B$ edge
use $\Sigma\alpha$, and collapse the $A$ edge.  On the face
use the barycentric coordinate of the common endpoint of the
long and $B$ edges as the suspension coordinate, with space
coordinate $\alpha(b)$.  On the third edge that coordinate
is an endpoint of the suspension, so the definitions agree.
This gives a continuous map
\[
 BC\longrightarrow\Sigma K.
\]
It induces an isomorphism on rational degree-two cohomology.
Indeed the triangle spectral sequence identifies that group
with
\[
 \ker\bigl(H^1(A)\oplus H^1(B)\oplus H^1(K)
        \longrightarrow H^1(A\times B)\bigr).
\]
The first two coordinates map isomorphically to the two
summands of $H^1(A\times B)$, so projection of this kernel
onto $H^1(K)$ is an isomorphism.  The suspension map just
constructed induces exactly its inverse, up to the harmless
edge orientation signs.  There is no additional degree-two
term: the map from the edge degree-zero groups to the face
degree-zero group is surjective.  This remains true if $K$
is disconnected.

A suspension has zero positive cohomology products, so all
degree-two products of $BC$ vanish in this case as well.
If $t\in V$, interchange the two adjacent edges.  These
possibilities exhaust $T$ and prove the theorem.
\end{proof}

\begin{remark}[The completed sixteen-arrow control]
\label{C:affine-chart16-control}
For $|T|=5$, \path{code/d_affine_chart16.py} enumerates all
$1536$ ordered pairs of retractions onto a common subset.
They give $24$ right-action orbits under state permutations.
Restoring every relative left-action position gives $3374$
commuting bisets and $12437$ equivariant composition choices.
Exactly $2311$ literal tables have $H_1(-;\mathbb F_2)=0$
and one-dimensional $H_2$.  The independent bar checker
verifies all $249588$ typed associativity equations and an
explicit zero-square primitive for every one of these tables.
The complete records have prefix
\path{code/outputs/d_affine_chart16_CP2}.
This calculation is a control of the universal theorem, not
its proof.
\end{remark}
